\begin{explanationbox}

	\textbf{\textcolor{CExplanation}{一句话直觉}}

	线面角的正弦值等于法向量与方向向量夹角余弦的绝对值，无需作辅助线，直接通过坐标计算得到。

	\medskip

	\textbf{\textcolor{CExplanation}{核心拆解}}

	\begin{description}[style=nextline, leftmargin=2.8em, labelsep=0.5em, itemsep=0.45em]
		\item[\textbf{要点一（几何关系）：}]
		      法向量 $\vec{n}$ 垂直于平面，方向向量 $\vec{l}$ 与 $\vec{n}$ 的夹角 $\varphi$ 满足 $\theta = |90^\circ - \varphi|$，因此 $\sin\theta = |\cos\varphi|$。

		\item[\textbf{要点二（代数转化）：}]
		      由 $\cos\varphi = \dfrac{\vec{n}\cdot\vec{l}}{|\vec{n}|\,|\vec{l}|}$，直接得到 $\sin\theta = \dfrac{|\vec{n}\cdot\vec{l}|}{|\vec{n}|\,|\vec{l}|}$。
	\end{description}

	\medskip

	\textbf{\textcolor{CExplanation}{几何本质（投影与垂直）}}

	直线与平面的夹角是其方向向量与平面内投影向量的夹角。由于法向量垂直于平面，方向向量在法向量上的投影分量与平面垂直，剩余分量与法向量垂直，可看作平面方向分量。因此线面角的正弦等于方向向量在法向量方向上的投影长度与方向向量长度之比。

	\medskip

	\textbf{\textcolor{CExplanation}{代数意义（点积比值）}}

	公式将角度计算转化为向量点积和模的比值，纯代数运算，适用于坐标已知或可计算的情形。结果与向量长度无关，只与方向有关。

	\medskip

	\textbf{\textcolor{CExplanation}{考点价值（建系直接计算）}}

	\begin{itemize}[leftmargin=*, itemsep=0.5em, topsep=0.4em]
		\item \textbf{考法一（已知法向量）：} 直接给出直线方向向量和平面法向量，或通过两点坐标求方向向量，代入公式得正弦值。
		\item \textbf{考法二（综合建系）：} 在立体几何证明题中，先建立空间直角坐标系，求出关键点坐标，进而得到方向向量和法向量，再套用公式。
	\end{itemize}

	\medskip

	\textbf{\textcolor{CExplanation}{顿悟点（法向量垂直关联）}}

	将线面角转化为法向量与方向向量的角，关键是对“垂直”的理解：法向量代表了平面的“朝向”，由 $\theta = |90^\circ - \varphi|$ 可知 $\sin\theta = |\cos\varphi|$。

	\medskip

	\textbf{\textcolor{CExplanation}{使用场景（向量法求角）}}

	\begin{itemize}[leftmargin=*, itemsep=0.5em, topsep=0.4em]
		\item \textbf{场景一（直接代入）：} 已知 $\vec{l}$ 和 $\vec{n}$ 的坐标，直接计算 $\dfrac{|\vec{l}\cdot\vec{n}|}{|\vec{l}|\,|\vec{n}|}$。
		\item \textbf{场景二（建系求解）：} 题目中包含垂直条件或可建立坐标系，先求方向向量和法向量，再代入公式。
	\end{itemize}

\end{explanationbox}